\SourceLeaf{118}
\par
Es ist also eins gegen eins zu wetten, dass der wahre Werth von \(h\)
\[
\text{zwischen}\quad H\bigl(1-\tfrac\rho{\surd m}\bigr)\quad\text{und}\quad H\bigl(1+\tfrac\rho{\surd m}\bigr)
\]
liegt, oder dass der wahre Werth von \(r\)
\[
\text{zwischen}\quad\frac{R}{1-\frac\rho{\surd m}}\quad\text{und}\quad\frac{R}{1+\frac\rho{\surd m}}
\]
falle, wenn wir durch \(R\) den im vorhergehenden Art. gefundenen wahrscheinlichsten Werth von \(r\) bezeichnen. Man kann diese Grenzen die \textit{wahrscheinlichen Grenzen der wahren Werthe von \(h\) und \(r\)} nennen; offenbar dürfen wir für die wahrscheinlichen Grenzen des wahren Werthes von \(r\) hier auch setzen
\[
R\bigl(1-\tfrac\rho{\surd m}\bigr)\quad\text{und}\quad R\bigl(1+\tfrac\rho{\surd m}\bigr)
\]

\Article{5}
Wir sind bei der vorhergehenden Untersuchung von dem Gesichtspunkte ausgegangen, dass wir \(\alpha,\,\beta,\,\gamma,\,\delta\) u.\,s.\,w. als bestimmte und gegebene Grössen betrachteten. und die Grösse der Wahrscheinlichkeit suchten, dass der wahre Werth von \(h\) oder \(r\) zwischen gewissen Grenzen liege. Man kann die Sache auch von einer andern Seite betrachten, und unter der Voraussetzung, dass die Beobachtungsfehler irgend einem bestimmten Wahrscheinlichkeitsgesetze unterworfen sind, die Wahrscheinlichkeit bestimmen, mit welcher erwartet werden kann, dass die Summe der Quadrate von \(m\) Beobachtungsfehlern zwischen gewisse Grenzen falle. Diese Aufgabe, unter der Bedingung, dass \(m\) eine grosse Zahl sei, ist bereits von \textsc{Laplace} aufgelöset, eben so wie diejenige, wo die Wahrscheinlichkeit gesucht wird, dass die Summe von \(m\) Beobachtungsfehlern selbst zwischen gewisse Grenzen falle. Man kann leicht diese Untersuchung noch mehr generalisiren; ich begnüge mich, hier das Resultat anzuzeigen.

Es bezeichne \(\varphi x\) die Wahrscheinlichkeit des Beobachtungsfehlers \(x\), so dass \(\int\varphi x.\mathrm{d}x=1\) wird, wenn man das Integral von \(x=-\infty\) bis \(x=+\infty\) ausdehnt. Zwischen denselben Grenzen wollen wir allgemein den Werth des Integrals
\[
\int\varphi x.x^n\mathrm{d}x
\]
\SourceLeaf{119}
durch \(K^{\mathrm n}\) bezeichnen\footnote{Oder vielmehr, das Integral \(\int\varphi x.x^n\mathrm{d}x\) zwischen den Grenzen \(x=0\) bis \(x=\infty\) soll durch
\[\tfrac12K^{\mathrm n}\]
bezeichnet werden. [Handschriftliche Bemerkung]}. Es sei ferner \(S^{\mathrm n}\) die Summe
\[
\alpha^n+\beta^n+\gamma^n+\delta^n+\text{ u.\,s.\,w.}
\]
wo \(\alpha,\,\beta,\,\gamma,\,\delta\) u.\,s.\,w. unbestimmt \(m\) Beobachtungsfehler bedeuten; die Theile jener Summe sollen, auch für ein ungerades \(n\), alle positiv genommen werden.

Sodann ist \(mK^{\mathrm n}\) der wahrscheinlichste Werth von \(S^{\mathrm n}\) und die Wahrscheinlichkeit, dass der wahre Werth von \(S^{\mathrm n}\) zwischen die Grenzen \(mK^{\mathrm n}-\lambda\) und \(mK^{\mathrm n}+\lambda\) falle,
\[
=\theta\frac{\lambda}{\surd[2m(K^{2\mathrm n}-K^{\mathrm n}K^{\mathrm n})]}
\]
Folglich sind die wahrscheinlichen Grenzen von \(S^{\mathrm n}\)
\[
mK^{\mathrm n}-\rho\surd[2m(K^{2\mathrm n}-K^{\mathrm n}K^{\mathrm n})]\quad\text{und}\quad
mK^{\mathrm n}+\rho\surd[2m(K^{2\mathrm n}-K^{\mathrm n}K^{\mathrm n})]
\]
Dieses Resultat gilt allgemein für jedes Gesetz der Beobachtungsfehler. Wenden wir es auf den Fall an, wo
\[
\varphi x=\frac h{\surd\pi}e^{-hhxx}
\]
gesetzt wird, so finden wir
\[
K^{\mathrm n}=\frac{\Pi\,\frac12(n-1)}{h^n\surd\pi}
\]
die Charakteristik \(\Pi\) in der Bedeutung der \textit{Disquisitiones generales circa seriem infinitam} (Comm. nov. soc. Gotting. T. II.) genommen (M. 5. Art. 28. der angef. Abh.) Also
\[
\begin{aligned}
K&=1,&K'&=\frac1{h\surd\pi},&K''&=\frac1{2hh},&K'''&=\frac1{h^3\surd\pi}\\[5pt]
K^{\mathrm{IV}}&=\frac{1.3}{4h^4},&K^{\mathrm{V}}&=\frac{1.2}{h^5\surd\pi},&K^{\mathrm{VI}}&=\frac{1.3.5}{8h^6},&K^{\mathrm{VII}}&=\frac{1.2.3}{h^7\surd\pi}\text{ u.\,s.\,w.}
\end{aligned}
\]
Es ist folglich der wahrscheinlichste Werth von \(S^{\mathrm n}\)
\[
\frac{m\Pi\,\frac12(n-1)}{h^n\surd\pi}
\]
und die wahrscheinlichen Grenzen des wahren Werthes von \(S^{\mathrm n}\)
